\section{Binary sections and joint character incidence}
\label{sec:binary-capacity}
The nonbinary hypothesis in Lemma~\ref{lem:unipotent} is essential.
For binary tops it is replaced by a joint budget indexed by actual
character kernels. Write $B_a=\binom a{\lfloor a/2\rfloor}$, including
$B_0=1$.

\begin{lemma}[Locally dependent maps over $\F_2$]\label{lem:locally-dependent}
Let $\mathcal M\leq\Hom(V,W)$ have dimension at least two, and suppose
that $\dim\mathcal M(v)\leq1$ for every $v\in V$. All maps in
$\mathcal M$ have image in one common line of $W$.
\end{lemma}
\begin{proof}
If $f\in\mathcal M$ has rank at least two, then for every
$g\in\mathcal M$ one has $\ker f\leq\ker g$. To see this, take
$u\in\ker f$ and two vectors with independent $f$-images. Applying
pointwise dependence to each vector and its translate by $u$ forces
$g(u)$ to lie in both independent lines. Thus $g$ induces a map on
$\operatorname{im}f$ taking each $x$ to either $0$ or $x$.
Linearity on independent pairs forces this map to be zero or the
identity. Hence $g=0$ or $f$, contradicting $\dim\mathcal M\geq2$.
All nonzero maps therefore have rank one. Two maps with distinct image
lines have the form $u\alpha$ and $v\beta$ for independent $u,v$.
There is an input on which both nonzero functionals $\alpha,\beta$
are one, contradicting pointwise dependence.
\end{proof}

\begin{lemma}[A deterministic separator]\label{lem:separator}
Let $\mathcal L\leq H\otimes S$ have dimension $\ell$, and put
\[
 j_0=\max_{0\ne\lambda\in H}
       \dim(\mathcal L\cap(\lambda\otimes S)),
\]
with empty maximum zero. There are $q$ independent functionals on $S$
whose simultaneous evaluation is injective on $\mathcal L$, where
\begin{equation}\label{eq:separator}
 q\leq\left\lfloor\frac{\ell+\max(1,j_0)}2\right\rfloor.
\end{equation}
For $\mathcal L=0$ take $q=0$.
\end{lemma}
\begin{proof}
View tensors as maps $S^*\to H$. On the current common kernel,
select an evaluation of rank at least two whenever one exists. Each
step removes at least two dimensions, and its functional is independent
of those already selected. If a remainder of dimension $r\geq2$
has only evaluations of rank at most one, Lemma~\ref{lem:locally-dependent}
puts it in one character line, so $r\leq j_0$. Separate its support
with $r$ more functionals. A remainder of dimension one requires only
one functional. If $k$ first-stage steps occurred, then
$\ell\geq2k+r$ and $q=k+r$, proving \eqref{eq:separator} in all cases.
The final functionals remain independent: earlier ones vanish on the
remaining support.
\end{proof}

Let $T\leq S_{2^a}$ be transitive and binary and $A$ a section of its
permutation module whose action factors through $B$. Set
\[
 S=A^B,\quad t=\dim S,\quad H=H^1(B,\F_2),\quad
 \mathcal L=\ker(H^1(B,S)\to H^1(B,A)),\quad\ell=\dim\mathcal L.
\]
The soluble module bound gives $t,\ell\leq B_a$, by applying it
to $A^B$ and the actual section $(A/S)^B\cong\mathcal L$.

\begin{lemma}[Joint character-kernel budget]\label{lem:joint-budget}
Inflate $H$ into $\Hom(T,C_2)$ and fix a point $x$. Let
$H_0=\{\chi\in H:\chi(T_x)=0\}$. For $W\leq H$ put
$j=\dim(W\cap H_0)$. Then
\begin{equation}\label{eq:joint-budget}
 t+\dim(\mathcal L\cap(W\otimes S))\leq2^jB_{a-j}.
\end{equation}
\end{lemma}
\begin{proof}
Let $K_W=\bigcap_{\chi\in W}\ker\chi$ in $T$. Its number of orbits
is $[T:K_WT_x]=2^j$, since the annihilator of the image of $T_x$ in
$T/K_W$ is $W\cap H_0$. The orbit lengths are $2^{a-j}$.
The preimage in $A$ of the connecting-map slice
$\mathcal L\cap(W\otimes S)$ has dimension equal to the left side
of \eqref{eq:joint-budget}. The group $K_W$ fixes it: the corresponding
cocycles into the trivial module $S$ vanish on $K_W$.
Filter this actual trivial section by the $K_W$-orbit modules and
apply the soluble bound $B_{a-j}$ to each. This proves the inequality.
\end{proof}

\begin{theorem}[Binary capacity from degree $64$]\label{thm:binary-large-capacity}
For $a\geq5$ there is $C\leq S$ with
\begin{equation}\label{eq:binary-capacity}
 2\dim C+4\dim(A/C)^B\leq2B_a+2B_{a-1}.
\end{equation}
For $a\geq6$ the right side is at most $15\cdot2^a/16$.
\end{theorem}
\begin{proof}
A nonzero character has one or two kernel orbits, so
Lemma~\ref{lem:joint-budget} gives
$j_0\leq J=2B_{a-1}-t$, using $B_a\leq2B_{a-1}$.
Apply Lemma~\ref{lem:separator} and let $C$ be the common kernel in
$S$ of the selected functionals. Then $\dim C=t-q$ and
$\mathcal L\cap(H\otimes C)=0$. Equation~\eqref{eq:fixed-quotient}
gives $\dim(A/C)^B=q$, so the cost is $2t+2q$.
If $J\geq1$, this is at most
$2t+\ell+J=t+\ell+2B_{a-1}\leq2B_a+2B_{a-1}$.
If $J=0$, then $t=2B_{a-1}\leq B_a$. The binomial identities imply
$a$ even and $t=B_a=2B_{a-1}$; this central binomial coefficient is
even. Therefore $q\leq\lceil\ell/2\rceil\leq B_a/2$, giving the
same result. Finally, for $a\geq6$, both central-binomial proportions
$B_a/2^a$ and $B_{a-1}/2^{a-1}$ are at most $5/16$.
\end{proof}

\begin{theorem}[The degree-$32$ boundary]\label{thm:binary32-capacity}
If $a=5$, there is $C\leq S$ with
$2\dim C+4\dim(A/C)^B\leq30$.
\end{theorem}
\begin{proof}
Here $t,\ell\leq10$ and $j_0\leq12-t$. The separator gives cost at
most $t+\ell+12\leq32$. Unless $t=\ell=10$, it is at most $31$
and hence at most $30$, since the cost is even.

Assume $t=\ell=10$. Put $H_0$ as in Lemma~\ref{lem:joint-budget}
and choose a complement $H=H_0\oplus H_1$. For
\[
 (M_0,M_1,M_2,M_3,M_4,M_5)=(10,12,12,16,16,32)
\]
the joint budget is $\dim(\mathcal L\cap(W\otimes S))\leq M_j-10$.
In particular $\mathcal L\cap(H_1\otimes S)=0$, so projection embeds
$\mathcal L$ into a ten-dimensional space
$\overline{\mathcal L}\leq H_0\otimes S$. For $W_0\leq H_0$,
\[
 \dim(\overline{\mathcal L}\cap(W_0\otimes S))
 \leq M_{\dim W_0}-10.
\]
We have $\dim H_0\leq5$ from the orbit count; taking $W_0=H_0$
forces equality, because $M_j-10\leq6$ for $j\leq4$.

Let $R_q$ be the number of nonzero tensors of rank $q$ in
$\overline{\mathcal L}$, $1\leq q\leq5$. Double count containment
of their row support in $r$-spaces of $H_0$. The instances $r=2,4$
give the simultaneous constraints
\begin{equation}\label{eq:rank-moments}
 \begin{split}
 15R_1+R_2&\leq465,\\
 15R_1+7R_2+3R_3+R_4&\leq1953,\\
 R_1+R_2+R_3+R_4+R_5&=1023.
 \end{split}
\end{equation}
Indeed the exact double count is
$\sum_{\dim W=r}(2^{\dim(\overline{\mathcal L}\cap(W\otimes S))}-1)
=\sum_qR_q\gauss{5-q}{r-q}2$.

For a uniform five-space $C\leq S=\F_2^{10}$, the probability of
containing a fixed $q$-dimensional image is
\[
 p_q=\frac{\gauss{10-q}{5-q}2}{\gauss{10}52},\qquad
 (p_1,\ldots,p_5)=
 \left(\frac1{33},\frac5{5621},\frac1{40953},
       \frac1{1733677},\frac1{109221651}\right).
\]
Set $z=p_5$, $y=(p_3-p_5)/3$, and $x=(p_1-15y-z)/15$.
These are positive, and the coefficient vector of $x$ times the first
row of \eqref{eq:rank-moments}, $y$ times the second, and $z$ times
the third dominates $(p_1,\ldots,p_5)$ coordinatewise. The slacks in
coordinates two and four are respectively
$1932416/1638324765$ and $2480/327664953$; the other three are equalities.
Consequently
\begin{equation}\label{eq:positive-dual}
 \sum_qp_qR_q\leq465x+1953y+1023z
 =\frac{103925485}{109221651}<1.
\end{equation}
Projection from $H$ to $H_0$ never increases tensor rank, and is
injective on $\mathcal L$. Since containment probability decreases
with rank, the expected number of nonzero tensors of the original
$\mathcal L$ contained in $H\otimes C$ is also less than one.
This number is an integer, so some $C$ has zero intersection.
Equation~\eqref{eq:fixed-quotient} then gives
$\dim(A/C)^B=10-5=5$, of cost $10+20=30$.
\end{proof}

Every simple binary module for a $2$-group is trivial. Hence these
costs are precisely $2\dim C+4r(A/C)$. A binary transitive pair
action of width $2s\geq64$ therefore has an entry certificate for
every normal axis, with $g\geq s/16$, by
Proposition~\ref{prop:prefix}. To sum over all widths one must also
bound the number of actions, axes, pair systems, cuts, and constants
$D_C$ uniformly. The finite-menu theorem alone does not perform that
unbounded sum.
